% 对应题号 M1-29；图与本题条件、推导、结论逐项对应。
% 图契约：印刷宽16cm；最小字号8.4pt；几何/场或曲线在左，计算关系在右。
\documentclass[tikz,border=7pt]{standalone}
\usepackage{fontspec,xeCJK,amsmath,amssymb,graphicx}
\IfFontExistsTF{Times New Roman}{\setmainfont{Times New Roman}}{\setmainfont{TeX Gyre Termes}}
\IfFontExistsTF{Heiti SC}{\setCJKmainfont[AutoFakeBold=1.5]{Heiti SC}}{\setCJKmainfont[AutoFakeBold=1.5]{FandolHei-Regular}}
\usetikzlibrary{arrows.meta,bending,calc,patterns,positioning}
\definecolor{ink}{HTML}{243745}
\definecolor{blue}{HTML}{275DAB}
\definecolor{green}{HTML}{177C63}
\definecolor{red}{HTML}{BD4935}
\definecolor{orange}{HTML}{AA6C16}
\definecolor{purple}{HTML}{715393}
\tikzset{every node/.style={font=\fontsize{9.3}{12.5}\selectfont,text=ink},
 axis/.style={-{Stealth[length=4pt]},line width=.65pt,draw=ink!65},
 vec/.style={-{Stealth[length=5pt]},line width=1pt,draw=blue},
 short/.style={-{Stealth[length=3pt]},line width=.8pt},
 guide/.style={dashed,draw=ink!40,line width=.5pt},
 card/.style={rounded corners=2pt,draw=blue!65,fill=blue!4,align=center,text width=6.6cm,inner xsep=9pt,inner ysep=7pt},
 title/.style={font=\fontsize{10.5}{14}\selectfont\bfseries,text width=15cm,align=center},
 label/.style={font=\fontsize{8.4}{11}\selectfont,align=center}}
\begin{document}
\begin{tikzpicture}[x=1cm,y=1cm]
\path[use as bounding box] (-8,-3.4) rectangle(8,3.4);
\node[title] at (0,3) {M1-29\quad 常矢量场穿入与穿出的通量抵消};
\begin{scope}[shift={(-4,0)}]
\draw[ink!55,fill=blue!3] (-1.2,-1.2) rectangle(1.2,1.2);
\draw[vec] (-2.1,-0.8)--(2.1,-0.8);
\draw[vec] (-2.1,0)--(2.1,0);
\draw[vec] (-2.1,0.8)--(2.1,0.8);

\draw[short,red] (-1.2,-1.55)--(-2,-1.55) node[left] {$-\mathbf n$};
\draw[short,red] (1.2,-1.55)--(2,-1.55) node[right] {$+\mathbf n$};
\node[label] at (0,-2.6) {两侧法向相反，场保持相同};

\end{scope}\node[card,draw=blue!65,fill=blue!4] at (3.95,1.65) {$\mathbf A=\mathbf A_0$ 为常矢量};
\node[card,draw=green!65,fill=green!4] at (3.95,0) {$\displaystyle\oint_{\partial V}\mathbf A_0\cdot d\mathbf S=\mathbf A_0\cdot\oint_{\partial V}d\mathbf S=0$};
\node[card,draw=orange!65,fill=orange!4] at (3.95,-1.65) {$\nabla\cdot\mathbf A_0=0$\\净通量为零，不要求每个面通量为零。};
\end{tikzpicture}\end{document}
